\scanpage{188}
\[
\Delta(x,y,z)^2=\det\begin{pmatrix}
\varphi(1,1)&\varphi(1,x)&\varphi(1,y)&\varphi(1,z)\\
\varphi(x,1)&\varphi(x,x)&\varphi(x,y)&\varphi(x,z)\\
\varphi(y,1)&\varphi(y,x)&\varphi(y,y)&\varphi(y,z)\\
\varphi(z,1)&\varphi(z,x)&\varphi(z,y)&\varphi(z,z)
\end{pmatrix}
\]
où $\varphi(u,v)$ est la forme bil. sym. associée à la forme quadratique $q$ ($=\det$)
\[
\varphi(u,v)=\operatorname{Tr}u\operatorname{Tr}v-\operatorname{Tr}(uv)=
\]
\uline{NB} On a
\[
\varphi(u,1)=\operatorname{Tr}(u)\quad\deleted{\unsure{\varphi(u,1)}}
\]
\[
\varphi(u,u)=2\det u
\]
Donc on trouve
\[
\boxed{\Delta(x,y,z)^2=\det\begin{pmatrix}
2&b&b'&b''\\
b&2c&bb'-t''&bb''-t'\\
b'&bb'-t''&2c'&b'b''-t\\
b''&bb''-t'&b'b''-t&2c''
\end{pmatrix}}
\]
ce qui permet d'expliciter
\[
4\tau\tau'=\Sigma^2-\Delta^2
\]
\margin{\unsure{où}}
où $\Delta$ est donné ci-dessus, et
\[
\boxed{\Sigma=bt+b't'+b''t''-bb'b''}
\]
\margin{187}
